\documentclass{article}
\usepackage{amsmath}
\usepackage{amssymb}
\usepackage{amsthm}
\usepackage{geometry}
\usepackage{graphicx}
\usepackage{hyperref}
\usepackage{natbib}
\geometry{margin=1in}

\title{Mind at Work: Detecting Agency through Assembly and Action}
\author{Flyxion}
\date{September 10, 2025}

\begin{document}

\maketitle

\begin{abstract}
This monograph presents a comprehensive framework for detecting mind-like agency by integrating the Effective Assembly Index for Mind Recognition (EAIMR) with the praxicon concept from cognitive science. EAIMR quantifies the improbability of a system’s structure arising from unguided processes, adjusted for efficiency and amplified by coherence, while the praxicon provides a structured repository of action representations bridging perceptual, semantic, and motor domains. By combining these, we propose a robust method to identify agentic intelligence across neural systems, collective organisms, ecosystems, and artificial intelligence, avoiding anthropocentric biases that marginalize non-narrative minds, such as those with aphantasia, anendophasia, or ecological intelligences. The framework is formalized using information theory, graph theory, category theory, and sheaf theory, with applications in neuroscience, AI alignment, ecological cognition, and cross-cultural studies. We address intersubjectivity collapse as a structural risk, amplified by Big Tech’s attention-compute arms race and ecological fragmentation, drawing parallels to historical critiques of industrialization. The monograph contrasts this approach with competing theories, provides detailed case studies, and discusses ethical implications for recognizing diverse minds.
\end{abstract}

\section{Introduction}
\label{sec:introduction}

The quest to detect and understand agency—defined as the capacity of a system to pursue goals through coordinated actions deviating from unguided physical processes—spans cognitive science, philosophy of mind, artificial intelligence (AI), and ecology. Our civilization’s bias toward equating intelligence with narrativity and imagery, dominant human cognitive modes, systematically marginalizes alternative forms of intelligence. As Michael Levin articulates, “mind blindness” leads us to dismiss the agency of plants, animals, and cellular collectives because they lack human-like symbols or narratives \citep{levin2022technological}. Similarly, cognitive minorities, such as individuals with aphantasia (inability to visualize mentally) or anendophasia (absence of inner monologue), are sidelined in an attention economy that rewards performative, narrative-driven cognition \citep{zeman2015, hurlburt2024}. This bias foreshadows intersubjectivity collapse (IC), where shared understanding among diverse minds breaks down, potentially leading to civilizational fragmentation \citep{pachniewski2023intersubjectivity}.

Big Tech’s business model—selling compute and storage while buying attention through algorithmic amplification—exacerbates this collapse. By spotlighting outlier success stories and suppressing systemic critiques, platforms create a narrative arms race that erodes trust and shared meaning \citep{zuboff2019age, wu2016attention}. This dynamic parallels historical ecological critiques, from Joni Mitchell’s “Big Yellow Taxi” lamenting the paving of paradise for parking lots to Jim Amos’s vision of a “Coruscant planet” where data centers and nuclear plants displace ecosystems for a compute-driven metaverse \citep{carr2010shallows}. Roads, like digital feeds, fragment habitats and meaning, isolating species and minds alike \citep{lanham2006economics}. The Effective Assembly Index for Mind Recognition (EAIMR) and the praxicon framework offer a solution by quantifying agency through structural complexity and action representation, bypassing narrativity bias to recognize diverse intelligences.

EAIMR measures the improbability of a system’s assembly, adjusted for reuse, parallelism, and degeneracy, and amplified by coherence \citep{walker2013algorithmic}. The praxicon, a cognitive “dictionary” of actions, bridges perceptual, semantic, and motor domains, grounding intelligence in embodied processes \citep{pastra2011praxicon, peigneux2004imaging}. Together, they detect agency in systems from neural circuits to termite mounds, forests, and AI, fostering inclusivity for non-human and atypical human minds. This monograph integrates these frameworks with information theory, graph theory, category theory, and sheaf theory, addressing applications in neuroscience, AI alignment, ecology, and cross-cultural cognition, while situating IC as a meta-crux for debates on AI futures, such as those between Yudkowsky’s pessimism and Miller’s cooperative optimism \citep{yudkowsky2008artificial, miller2006robust}.

\section{Mind Blindness and Intersubjectivity Collapse}
\label{sec:mind_blindness}

Mind blindness, as described by Levin, is the failure to recognize intelligence in systems lacking human-like symbolic or narrative expression \citep{levin2022technological}. Intelligence, in this view, is any system’s ability to pursue goals, adapt to perturbations, and solve problems in its medium, from regenerating limbs to forest nutrient cycling \citep{levin2019computational}. This bias extends to cognitive minorities—individuals with aphantasia (2–4\% of the population, unable to visualize mentally) or anendophasia (lacking inner speech)—who face exclusion in narrative-driven economies despite intact functional reasoning \citep{jacobs2022, bps_anendophasia}. Intersubjectivity collapse, as framed by Pachniewski, occurs when radically different minds render shared norms and trust unintelligible, a risk amplified by AI, cognitive augmentation, and platform dynamics \citep{pachniewski2023intersubjectivity}.

\subsection{Big Tech and the Arms Race of Intersubjectivity}
Big Tech’s business model—selling compute (AWS, GCP, Azure) and buying attention (via social feeds and ads)—creates an arms race of narratives. Platforms amplify rare successes (e.g., a single coding bootcamp graduate’s story) while burying critiques, skewing perceptions of efficacy \citep{zuboff2019age, wu2016attention}. This mirrors Mark Miller’s arms race logic, where destructive strategies (misleading narratives) outpace defensive ones (criticism, regulation) due to scalability \citep{miller1988markets}. Programming courses exemplify this: an oversupply of training with undersupplied job demand creates a hope market, where platforms profit by spotlighting outliers \citep{carr2010shallows}. This erodes trust, fostering cynicism or factional belief, a proto-collapse of shared meaning.

Ecologically, this parallels the critique of industrialization. Joni Mitchell’s “Big Yellow Taxi” mourned paved-over ecosystems; Jim Amos’s “Coruscant planet” envisions data centers displacing water and air for compute \citep{lanham2006economics}. Roads fragment habitats—disrupting migrations, raising heat islands, and diverting water—much as digital feeds fragment intersubjectivity into echo chambers \citep{arendt1958human}. Both create isolated patches, undermining exchange of animals or ideas, driving collapse.

\subsection{Mind Blindness Across Domains}
1. **Ecological Minds (Plants and Animals)**  
   Bias: Intelligence requires language or narrative.  
   Blindness: Plants and animals use chemical gradients, root networks, and collective behaviors, dismissed as “instinct” \citep{gieling2011assessing, mendell2010why}.  
   Consequence: Ecological exploitation due to unrecognized agency.

2. **Cognitive Minorities (Aphantasia, Anendophasia, Atypical Cognition)**  
   Bias: Value lies in attention-driven narration and imagery.  
   Blindness: Aphantasic and anendophasic individuals process functionally but struggle in performative economies \citep{zeman2015, hurlburt2024}.  
   Consequence: Social and economic exclusion in platform-optimized systems.

3. **Religious and Philosophical Systems**  
   Bias: Each tradition’s narrative (e.g., Bible, Qur’an, scientific method) is universal.  
   Blindness: Rival systems’ values become illegible, leading to mistrust \citep{habermas1984theory}.  
   Consequence: Cultural fragmentation and conflict.

Unifying Thread: Mind blindness drives intersubjectivity collapse by excluding intelligences that don’t conform to narrative norms, from ecosystems to cognitive minorities to AI \citep{cabrera2013human}.

\subsection{Termite–Neuron–Forest Triad}
The Termite–Neuron–Forest Triad illustrates emergent intelligence from swarms of simpler agents, bypassing narrativity bias \citep{levin2022technological}.
1. **Neurons as “Just Insects”**: Synaptic swarms yield consciousness, with high assembly indices reflecting improbable organization \citep{gallese1998mirror}.
2. **Termites as “Just Insects”**: Pheromone-driven swarms build climate-controlled mounds, with assembly indices exceeding random soil structures \citep{turner2000extended}.
3. **Trees as “Just Plants”**: Mycorrhizal networks create resilient ecosystems, with high assembly indices signaling collective agency \citep{camazine2001self}.

EAIMR quantifies these as \(EA(X) = A(X) / (R \cdot P \cdot D)\), amplifying coherence to detect mind-like processes \citep{walker2013algorithmic}.

\section{Historical and Conceptual Foundations}
\label{sec:foundations}

\subsection{Prerequisites: Complexity and Information}
Complexity distinguishes agentic systems from inert matter. Murray Gell-Mann’s effective complexity highlights balanced regularities \citep{adami2012use}. Hazen’s functional information measures improbable configurations serving roles, foundational to assembly theory \citep{hazen2007functional}. Shannon’s entropy and Kolmogorov complexity provide quantitative tools for improbability \citep{shannon1948mathematical, kolmogorov1965three}.

\subsection{Assembly Theory and Complexity}
Assembly theory quantifies the minimal steps to construct structures, with high indices indicating biotic or agentic processes \citep{cronin2016beyond, walker2013algorithmic]. EAIMR extends this with adjustments for reuse, parallelism, degeneracy, and coherence, addressing loopholes like scaffolded complexity \citep{adami2012use}.

\subsection{Praxicon: Action Representation}
The praxicon, from apraxia studies, distinguishes input (recognition) and output (production) action representations, grounded in parietal-premotor circuits \citep{peigneux2004imaging, mantovani2010reviewing}. It addresses the symbol grounding problem in AI, linking language to motor primitives \citep{harnad1990symbol, pastra2011praxicon, vannuscorps2023effector}.

\subsection{Competing Approaches}
EAIMR-praxicon contrasts with Integrated Information Theory (IIT) \citep{tononi2004information}, Free Energy Principle (FEP) \citep{friston2010free}, and enactive cognition \citep{barandiaran2009defining}, offering computable metrics for distributed agency.

\section{Formal Definitions}
\label{sec:formal}

\subsection{Prerequisites: Generative Models}
Null models \(P_0(X)\) represent mindless processes; mind-like models \(P_1(X)\) incorporate agency \citep{cover2006elements}. These underpin likelihood and surprise calculations.

\subsection{EAIMR Definition}
For system \(X\) over alphabet \(\Sigma\):

\[
\text{EAIMR}(X) = \frac{A(X)}{R(X) \cdot P(X) \cdot D(X)} \cdot C(X),
\]

where \(A(X)\), \(R(X)\), \(P(X)\), \(D(X)\), and \(C(X)\) are assembly index, reuse, parallelism, degeneracy, and coherence \citep{walker2013algorithmic].

\subsection{Praxicon Definition}
The praxicon is a bipartite graph \(G_{\mathcal{P}} = (V_A, V_S, E)\), linking action vertices \(V_A\) to semantic vertices \(V_S\) \citep{panunzi2018imagact, barsalou2008grounded}.

\section{Mathematical Integration}
\label{sec:integration}

\subsection{Information-Theoretic Foundations}
Self-information: \(I_0(X) = -\log P_0(X)\). Likelihood ratio: \(\Lambda(X) = P_1(X)/P_0(X)\). Action entropy reduction: \(\Delta H_{\mathcal{P}} = H(\mathcal{P}_0) - H_{\text{obs}}(\mathcal{P})\), with \(C(X) \approx 1 + \alpha \Delta H_{\mathcal{P}} / H(\mathcal{P}_0)\) \citep{shannon1948mathematical, jaynes2003probability}. Kolmogorov complexity approximates \(A(X)\) \citep{kolmogorov1965three].

\subsection{Graph-Theoretic Formalism}
Praxicon graph \(G_{\mathcal{P}}\) uses spectral properties for coherence \citep{newell1990unified]. EAIMR identifies low-entropy subgraphs.

\subsection{Category-Theoretic Interpretation}
The praxicon is a functor \(F: \mathcal{C}_{\text{perc}} \to \mathcal{C}_{\text{motor}}\) via \(\mathcal{C}_{\text{sem}}\). EAIMR measures natural transformation costs \(\eta: F \Rightarrow G\), reflecting improbable mappings \citep{healy2022, ehresmann2021}.

\subsection{Sheaf-Theoretic Interpretation}
A sheaf \(\mathcal{F}\) on system components assigns local praxicon fragments, with cohomology \(H^k(X, \mathcal{F})\) indicating non-trivial agency \citep{spivak2024, stenzel2024, ehresmann1980}.

\section{Worked Examples}
\label{sec:examples}

\subsection{Neural Tissue}
High \(A(X)\), praxicon engrams reduce entropy. Mirror neuron data \citep{gallese1998mirror]. Entropy curve: Sharp reduction.

\subsection{Termite Mound}
Moderate \(A(X)\), high coherence in thermoregulation \citep{turner2000extended, camazine2001self}. Graph: Worker actions to colony semantics.

\subsection{Forest Ecosystem}
High \(A(X)\), praxicon-like nutrient routing \citep{levin2019computational]. Lattice: Trophic interactions.

\subsection{Negative Case: Crystal Growth}
Low \(C(X)\), high predictability \citep{hazen2007functional].

\section{Applications}
\label{sec:applications}

\subsection{Neuroscience}
Detects apraxia and atypical cognition \citep{peigneux2004imaging, zeman2015].

\subsection{AI and Robotics}
Grounds LLMs; detects agency \citep{pastra2011praxicon, anderson2007how].

### Ecological and Collective Minds}
Identifies swarm intelligence \citep{camazine2001self, simondon2017mode].

### Cross-Cultural Cognition}
Rituals as praxicons \citep{ortega1961meditations].

\section{Limits and Future Research}
\label{sec:limits}

Challenges: Approximating complexity, validating sheaves. Future: Bioelectric models, AI simulations \citep{levin2019computational].

\section{Ethical and Philosophical Implications}
\label{sec:ethics}

Mitigates mind blindness, informs AI ethics, and aligns with process philosophy \citep{dennett1987intentional, whitehead1929process, deleuze1987thousand].

\section{Decision Rule}
\label{sec:decision}

\(\text{EAIMR}(X) \geq \tau\) classifies agency, using graph entropy proxies.

\section{Conclusion}
\label{sec:conclusion}

This framework detects agency, averting intersubjectivity collapse by recognizing diverse minds.

\appendix

\section{Mathematical Appendix: EAIMR, Information, and Improbable Order}
\label{app:math}

\subsection{Preliminaries \& Notation}
Let \(X\) denote a structure/system (e.g., termite mound, forest network, neural microcircuit) described over an alphabet \(\Sigma\).

Let \(P_0\) be a null (mindless) generative model for \(X\) (e.g., stochastic physics + simple self-organization) with distribution \(P_0(X)\).

Let \(P_1\) be a mind-like generative model for \(X\) (agentic/goal-directed) with distribution \(P_1(X)\).

Raw Assembly Index \(A(X)\): minimum step count in a prescribed assembly grammar to build \(X\) from primitives.

EAIMR:
\[
\boxed{\text{EAIMR}(X)=\frac{A(X)}{R(X)\,P(X)\,D(X)}\; C(X)}
\]

\subsection{Connection to Kolmogorov Complexity}
Let \(K(X)\) be prefix Kolmogorov complexity (shortest universal description of \(X\)) \citep{kolmogorov1965three].

In general, lower \(K(X)\) implies higher compressibility (simple regularities).

Many mind-built artifacts have mixed profiles: globally compressible high-level plan + low-level details that are costly to specify.

A useful proxy is two-part MDL (minimum description length):

\[L(X) \approx L(\text{model}) + L(\text{residuals}\mid\text{model})\]

Heuristic link (non-rigorous but actionable):

\[\text{EAIMR}(X) \;\propto\; \frac{K_\text{plan}(X)}{R\,P\,D} \cdot C \quad\text{with}\quad K_\text{plan}(X)\lesssim K(X)\]

\subsection{Statistical Surprise \& Likelihood Ratios}
Define the self-information under the null:

\[I_{0}(X) \equiv -\log P_{0}(X)\]

Likelihood ratio:

\[\Lambda(X) \equiv \frac{P_{1}(X)}{P_{0}(X)} \;=\; \exp\Big(\log P_{1}(X) - \log P_{0}(X)\Big).\]

High EAIMR correlates with large \(\Lambda(X)\) and high \(I_0(X)\), indicating \(X\) is unlikely under mindless dynamics but plausible under mind-like dynamics \citep{shannon1948mathematical].

A practical surrogate score:

\[S(X) \;\equiv\; I_{0}(X)\;+\;\lambda\, \phi(X)\]

\[\text{EAIMR}(X) \;=\; g\big(S(X)\big),\quad g\uparrow\]

\subsection{Entropy Reduction, Work, and Directed Assembly}
Let \(\mathcal{M}\) be macrostate features (e.g., ventilation topology, nutrient routing, neural coding). Define macro-entropy under \(P_0\) and observed for \(X\).

Directed assembly manifests as entropy reduction:

\[\Delta H \;\equiv\; H(\mathcal{M}) - H_\text{obs}(\mathcal{M}) \;\;>\; 0\]

A convenient coherence multiplier:

\[C(X) \;\approx\; 1 + \alpha \,\frac{\Delta H}{H(\mathcal{M})}\]

\subsection{Reuse, Parallelism, Degeneracy as Information Discounts}
Treat these as code-length discounts relative to the raw assembly plan \citep{hazen2007functional].

\textbf{Reuse}: Library motifs shorten description. If motif dictionary has description amortized across instances, effective bits drop. Model as a divisor on \(A(X)\).

\textbf{Parallelism}: Independent sub-assemblies compress as a product code; wall-clock and effective complexity per unit time fall. Model as a divisor on \(A(X)\).

\textbf{Degeneracy}: Many alternatives achieve the same function, reducing description length. Model as a divisor.

This aligns EAIMR with MDL-style accounting: EAIMR rises when coordinated specification remains large and functionally tight.

\subsection{Coherence via Control \& Causal Graphs}
Let \(G(X)\) be a causal graph extracted from \(X\)’s dynamics. Define a control coherence term:

\[\phi(X) \;=\; \beta_1\, \text{Stab}(G) \;+\; \beta_2\, \text{GoalAcc}(G) \;+\; \beta_3\, \text{Robust}(G)\]

\(\text{GoalAcc}(G)\): accuracy tracking a specified functional setpoint (e.g., temperature band in termite mounds).

\(\text{Robust}(G)\): performance under noise/lesions (graceful degradation).

Use \(\phi(X)\) in surrogate score or as multiplier.

\subsection{Practical Estimation (Finite Data)}
Because \(K(X)\) is uncomputable, rely on approximations:

\textbf{Compression proxies}: On suitable encodings (graphs, fields, symbol streams).

\textbf{Surprisal under null}: Fit \(P_0\) (e.g., generative physics sims, random graphs) and estimate by Monte Carlo or density models.

\textbf{Function \& coherence}: Quantify with control/evaluation tasks (tracking error, energy efficiency, resilience curves).

\textbf{Discounts}: Measure from libraries (motif counts), concurrency (critical path vs. total ops), and function-equivalence classes (number of viable variants).

\subsection{Worked Intuitions (Triad Examples)}
\textbf{Neural tissue}: Very high \(A(X)\) (precise mesoscale organization), strong \(C(X)\) (robust goal pursuit), discounts present (columns, parallelism), yet EAIMR stays high mind-like \citep{gallese1998mirror].

\textbf{Termite mound}: Moderate \(A(X)\) vs. dunes, high \(C(X)\) (thermoregulation), sizable discounts (motifs, parallel labor, many workable designs), but coordinated control keeps EAIMR high \citep{turner2000extended].

\textbf{Forest}: Extremely high macro \(A(X)\) vs. random vegetation, high robustness/goal achievement (nutrient routing, succession), massive reuse/degeneracy, yet coherence over centuries pushes EAIMR high \citep{levin2019computational].

\subsection{Decision Rule (Operational)}
Define thresholds from a reference set of clearly mindless assemblies:

\[\text{EAIMR}(X) \;\ge\; \tau \implies \text{classify as mind-like (candidate)}\]

\section{Derivations and Pseudocode}
\label{app:derivations}

\subsection{EAIMR Derivation}
From surprisal: \(S(X) = I_0(X) + \lambda \phi_{\mathcal{P}}(X)\), EAIMR as monotonic \(g(S(X))\).

\subsection{Pseudocode}
\begin{verbatim}
def compute_eaimr(X, P0, P1, G_P):
    A = assembly_index(X)
    R = reuse_factor(X)
    P = parallelism(X)
    D = degeneracy(X)
    C = coherence(G_P)
    return (A / (R * P * D)) * C
\end{verbatim}

\bibliographystyle{plainnat}
\bibliography{reference}

\end{document}